\subsection{八元数}

设 E 是 A 上类型为 \((\alpha, \beta, \gamma)\) 的四元数代数（§ 2，第 5 号，例 2），并设 \(\delta \in A\) 。E 的由 \(\delta\) 与 E 的共轭所定义的凯莱扩张 F 称为 A 上的八元数代数，并称其类型为 \((\alpha, \beta, \gamma, \delta)\) 。由第 2 号的命题 3，F 是交错代数。它有一个由 8 个元素组成的基 \((e_{i})_{0 \leq i \leq 7}\) ，定义为

\[
\begin{array}{l l l l} e _ {0} = (e, 0), & e _ {1} = (i, 0), & e _ {2} = (j, 0), & e _ {3} = (k, 0) \\ e _ {4} = (0, e), & e _ {5} = (0, i), & e _ {6} = (0, j), & e _ {7} = (0, k) \end{array}
\]

（其中 \((e, i, j, k)\) 是 E 的基，定义于 loc. cit.；显然 \(e_0\) （也记作 \(e\) ）是 F 的单位元。若 \(u = \sum_{i=0}^{7} \xi_i e_i\) 是 F 的一个元素（其中 \(\xi_i \in A\) ），则 § 2 第 5 号的公式 (23)、(24) 和 (31) 给出八元数 \(u\) 的共轭、迹和范数

\[
\left\{ \begin{array}{c} \bar {u} = (\xi_ {0} + \beta \xi_ {1}) e _ {0} - \sum_ {i = 1} ^ {7} \xi_ {i} e _ {i} \\ \mathrm{T} _ {\mathrm{F}} (u) = 2 \xi_ {0} + \beta \xi_ {1} \\ \mathrm{N} _ {\mathrm{F}} (u) = \xi_ {0} ^ {2} + \beta \xi_ {0} \xi_ {1} - \alpha \xi_ {1} ^ {2} - \gamma (\xi_ {2} ^ {2} + \beta \xi_ {2} \xi_ {3} - \alpha \xi_ {3} ^ {2}) \\ - \delta (\xi_ {4} ^ {2} + \beta \xi_ {4} \xi_ {5} - \alpha \xi_ {5} ^ {2}) + \gamma \delta (\xi_ {6} ^ {2} + \beta \xi_ {6} \xi_ {7} - \alpha \xi_ {7} ^ {2}). \end{array} \right.\tag{4}
\]

现在设 \(u = (x, y)\) , \(u' = (x', y')\) 和 \(u'' = (x'', y'')\) 是三个八元数（其中元素 \(x, x', x'', y, y', y''\) 属于 E）。§ 2 第 5 号的公式 (24) 和 (27) 给出

\[
\begin{array}{r l} & {\mathrm{T} _ {\mathrm{F}} ((u u ^ {\prime}) u ^ {\prime \prime}) = \mathrm{T} (x x ^ {\prime} x ^ {\prime \prime}) + \delta \mathrm{T} (\bar {y} ^ {\prime} y x ^ {\prime \prime}) + \delta \mathrm{T} (\bar {y} ^ {\prime \prime} y \bar {x} ^ {\prime}) + \delta \mathrm{T} (\bar {y} ^ {\prime \prime} y ^ {\prime} x)} \\ & {\mathrm{T} _ {\mathrm{F}} (u (u ^ {\prime} u ^ {\prime \prime})) = \mathrm{T} (x x ^ {\prime} x ^ {\prime \prime}) + \delta \mathrm{T} (x ^ {\prime \prime} \bar {y} ^ {\prime} y) + \delta \mathrm{T} (\bar {x} ^ {\prime} \bar {y} ^ {\prime \prime} y) + \delta \mathrm{T} (x \bar {y} ^ {\prime \prime} y ^ {\prime})} \end{array}
\]

（其中 T 表示 E 中的迹，并利用了 E 是结合的这一事实）。由于 \(\mathrm{T}(xy) = \mathrm{T}(yx)\) 对所有四元数 x、y 成立（§ 2，第 4 号，公式 (17)），由此可得

\[
\mathrm{T} _ {\mathrm{F}} ((u u ^ {\prime}) u ^ {\prime \prime}) = \mathrm{T} _ {\mathrm{F}} (u (u ^ {\prime} u ^ {\prime \prime})).\tag{5}
\]

我们特别研究类型为 \((-1, 0, -1, -1)\) 的八元数；此时公式（4）简化为

\[
\left\{ \begin{array}{c} {\bar {u}} = \xi_ {0} e _ {0} - \sum_ {i = 1} ^ {7} \xi_ {i} e _ {i} \\ \mathrm{T} _ {\mathrm{F}} (u) = 2 \xi_ {0} \\ \mathrm{N} _ {\mathrm{F}} (u) = \sum_ {i = 0} ^ {7} \xi_ {i} ^ {2}. \end{array} \right.\tag{6}
\]

*如果我们取 A 为实数域 R，则 R 上类型为 \((-1,0,-1,-1)\) 的八元数称为凯莱八元数（或 octaves）。由第 2 号的命题 2 (ii) 可知，每个 \(\neq0\) 的凯莱八元数都是可逆的。*

命题 4. 设 F 是 A 上类型为 \((-1,0,-1,-1)\) 的一个八元数代数。存在一个以二元素域 Z/2Z 为标量域的三维向量空间 V，以及一个双射 \(\lambda\mapsto e_{\lambda}^{\prime}\) （从 V 到基 \((e_{i})_{0\leq i\leq7}\) 上的双射），使得

\[
e _ {0} ^ {\prime} = e _ {0}, \quad e _ {\lambda} ^ {\prime} e _ {\mu} ^ {\prime} = \pm e _ {\lambda + \mu} ^ {\prime}\tag{7}
\]

对所有 \(\lambda, \mu\) （属于 V）。要使 \(e_{\lambda}'(e_{\mu}'e_{\nu}') = (e_{\lambda}'e_{\mu}')e_{\nu}'\) ，只需 \(\lambda, \mu, \nu\) （ V 中的）在 Z/2Z 上线性相关即可；若 \(2 \neq 0\) （在 A 中），则此条件是必要的。

我们保留本节开头的记号。由 § 2 第 5 号的公式 (33) 可知，由元素 \(\pm e_0\) , \(\pm e_1\) , \(\pm e_2\) , \(\pm e_3\) 组成的集合 S 在乘法下是稳定的。此外，对 \(x, y, y'\) （属于 E），由 § 2 第 5 号的公式 (22)，

\[
(x, 0) (0, y ^ {\prime}) = (0, y ^ {\prime} x), \quad (0, y ^ {\prime}) (x, 0) = (0, y ^ {\prime} \bar {x}), \quad (0, y) (0, y ^ {\prime}) = (- \bar {y} ^ {\prime} y, 0)\tag{8}
\]

使得由元素 \(\pm e_{i}\) ( \(0 \leqslant i \leqslant 7\) ) 组成的集合 T 在乘法下是稳定的；此外，其乘法表与环 A 无关。

特别地，设 \(\mathbf{A}''\) 为含两个元素的域 \(\mathbf{Z}/2\mathbf{Z}\) ，设 \(\mathbf{E}''\) 为 \(\mathbf{A}''\) 上类型为 (1, 0, 1) 的四元数代数， \(\mathbf{F}''\) 为 \(\mathbf{A}''\) 上类型为 (1, 0, 1, 1) 的八元数代数；令 \((e_i'')_{0 \leq i \leq 7}\) 为 \(\mathbf{F}''\) 的按上述方式构成的基。由于 \(-e_i'' = e_i''\) ，集合 \(\mathbf{T}''\) （由 \(e_i''\) 组成）有 8 个元素且在乘法下封闭；此外，由上述内容立即得出，映射 \(\theta: \mathbf{T} \to \mathbf{T}''\) （满足 \(\theta(e_i) = \theta(-e_i) = e_i''\) ，对所有 \(0 \leq i \leq 7\) ）是乘法同态。此外，四元数代数 \(\mathbf{E}''\) 在此情形下是交换的，因而 \(\mathbf{F}''\) 是结合的（ \(\S 2\) ，第 5 号，命题 5）；进一步， \(\mathbf{F}''\) 中的共轭在此情形下是恒等映射。因此 \(\mathbf{T}''\) 是一个群，且公式 (8) 表明它是交换的；这些公式以及 \(\S 2\) 第 5 号的公式 (33) 表明 \(\mathbf{T}''\) 的每个元素的平方都是单位元。若用 V 表示写成加法群后的群 \(\mathbf{T}''\) ，则可在 V 上赋予唯一的向量空间结构，其标量域为 \(\mathbf{Z}/2\mathbf{Z}\) ，其维数必然为 3，因为

\[
\operatorname{Card} (\mathbf {V}) = 8 = 2 ^ {\dim (\mathbf {V})}.
\]

对所有 \(\lambda \in V\) ，随后设 \(e_{\lambda}'\) 表示 \((e_i)_{0 \leq i \leq 7}\) 中满足 \(\theta(e_{\lambda}') = \dot{\lambda}\) 的元素；于是 \(e_0' = e_0\) ；此外，由于 \(\theta\) 是同态，且关系 \(\theta(x) = \theta(y)\) 等价于 \(x = \pm y\) ，故有 \(e_{\lambda}' e_{\mu}' = \pm e_{\lambda + \mu}'\) 。若 \(\lambda, \mu, \nu\) 在 \(\mathbf{Z}/2\mathbf{Z}\) 上线性无关，它们就构成 \(V\) 的一组基，从而所有元素 \(e_i\) ( \(0 \leqslant i \leqslant 7\) ) 都将属于由 \(e_{\lambda}', e_{\mu}'\) 和 \(e_{\nu}'\) 生成的子代数；当 \(2 \neq 0\) （在 \(A\) 中）时， \(e_{\lambda}'(e_{\mu}' e_{\nu}') = (e_{\lambda}' e_{\mu}') e_{\nu}'\) 是不可能的，因为那样 \(F\) 就是结合的，从而 \(E\) 是交换的（ \(\S 2\) ，第 5 号，命题 5），而这与 \(\S 2\) 第 5 号的关系式 (33) 相矛盾。另一方面，如果 \(\lambda, \mu, \nu\) 在 V 中线性

相关，则这三个元素 \(e_{\lambda}^{\prime}\) , \(e_{\mu}^{\prime}\) , \(e_{\nu}^{\prime}\) 属于 F 的一个由 2 个生成元生成的子代数，因而该子代数是结合的（第 1 号，命题 1）；结论由此得出。

注记。由于 \(\bar{e}_{\lambda}^{\prime} = -e_{\lambda}^{\prime}\) （对 \(\lambda \neq 0\) ），

\[
e _ {\lambda} ^ {\prime 2} = - e _ {0} \quad \text { for } \lambda \neq 0,
\]

\[
e _ {\mu} ^ {\prime} e _ {\lambda} ^ {\prime} = - e _ {\lambda} ^ {\prime} e _ {\mu} ^ {\prime} \quad \text { for } \lambda \neq 0, \mu \neq 0 \text { and } \mu \neq \lambda .
\]

